\chapter{Conclusions}
\label{sec:conclusions}

This thesis explored the use of deep learning as a unifying methodological framework for the study of physical systems in quantum information. Two distinct but conceptually related problems were addressed: the interpolation of unitary time-evolution operators generated by time-dependent Hamiltonians (Chapter~\ref{ch:results-pinn-unitaries}), and the reconstruction of global density matrices compatible with a prescribed set of quantum marginals (Chapter~\ref{ch:results-qmp}). In both cases, the central idea was the same: rather than treating neural networks as opaque function approximators, we embedded the known physical structure of each problem directly into the architecture and the training objective, so that the learned solutions remained consistent with the underlying quantum mechanics. This physics-informed perspective, introduced in Chapter~\ref{sec:deep_learning} and developed throughout the applications, is what gives both methods their accuracy, data efficiency, and physical interpretability.

\section{Summary of the main results}
\label{sec:conc-summary}

The first application demonstrated that a Physics-Informed Neural Network can learn the full time evolution of a closed quantum system from a sparse set of sampled unitaries, interpolating the time-evolution operator $U(t)$ across the entire time domain without ever being given the Hamiltonian explicitly. By incorporating unitarity as a soft constraint and, for the larger systems, by predicting an effective Hamiltonian propagated through a second-order Magnus expansion or a Trotterization, the framework achieved mean gate fidelities close to $0.99$ with low variability across independent Hamiltonian realizations, for systems ranging from 2 to 8 qubits. Remarkably, models trained with coarse temporal sampling ($\Delta t = 1.0$) reproduced the dynamics almost as well as those trained with fine grids, confirming a strong interpolation capability that translates directly into a reduction of the measurement cost associated with quantum process tomography. Beyond accuracy, the trained models offered a substantial computational advantage at inference time, with speedups of roughly one order of magnitude over a geodesic-interpolation baseline~\cite{schilling2024} that scales as $\mathcal{O}(2^{3N})$. This advantage is expected to grow with system size, making the approach particularly attractive in settings where the propagator must be evaluated repeatedly, such as dynamical quantum state tomography~\cite{Siva2022, peruzzo2026}.

The second application showed that the combination of a convolutional denoising autoencoder with the Marginal Imposition Operator~\cite{uzcategui2022mio} can construct physically valid global density matrices compatible with a given set of quantum marginals, a specific instance of the Quantum Marginal Problem, which is known to be QMA-complete in its general form~\cite{liu2006consistency, broadbent2022qma}. By encoding density matrices as two-channel images and designing a loss function that enforces Hermiticity, positivity, and normalization, the model preserved the target marginals with high fidelity for systems of 3 to 8 qubits. The hybrid \textit{model1}+MIO scheme reconstructed states whose marginals matched the targets with perfect fidelity in many cases, and once trained, the inference was significantly faster than state-of-the-art semidefinite programming solvers, even succeeding in regimes where the SDP solver failed to run. The successful use of transfer learning across system sizes further suggested that the network internalizes structural features of the compatibility problem that are not tied to a fixed Hilbert-space dimension.

Taken together, these two results illustrate a common and powerful pattern. In both cases, a quantum-information task that is computationally demanding by conventional means, integrating a non-commuting time-dependent Hamiltonian or searching the exponentially large space of compatible global states, was recast as a learning problem in which physical constraints restrict the hypothesis space to physically admissible solutions. The neural networks did not merely fit data; they learned representations that respect unitarity, positivity, and the partial-trace structure of composite systems. This is precisely the promise of physics-informed learning~\cite{Raissi2019, karniadakis2021}: combining the flexibility and scalability of deep learning with the rigor of the governing physical laws.

\section{Toward large-scale deep learning for quantum information}
\label{sec:conc-scale}

A recurring theme throughout this thesis is the exponential growth of the Hilbert space with the number of qubits, which renders exact classical methods intractable and motivates the search for scalable approximations~\cite{lloyd1996}. Both frameworks developed here confront this challenge directly, and both point toward a broader role for deep learning in the study of large-scale quantum systems.

The interpolation framework addressed scalability by shifting, for the largest systems, from learning the full $2^N \times 2^N$ unitary to learning a compact effective Hamiltonian with only $\mathcal{O}(10^3)$ parameters, propagated through structure-preserving numerical integrators. This decoupling of the learned representation from the dimension of the operator it generates is a general design principle: the network learns the few physically relevant degrees of freedom, while the exponential structure is restored by a physics layer that guarantees unitarity by construction. The marginal-reconstruction framework, in turn, exploited the locality-based feature extraction of convolutional architectures and the transferability of learned representations across system sizes, suggesting that a single architecture can adapt to growing Hilbert spaces with limited retraining rather than full redesign.

These observations support a more general claim. Deep learning is not merely a faster numerical solver for isolated quantum-information tasks; it is a framework for learning the latent, low-dimensional structure that often underlies high-dimensional quantum problems. Once that structure is captured, predictions can be generated at arbitrary points, such as at any time in the evolution or for any rank of the global state, at a computational cost that grows far more slowly than that of the exact methods they approximate. As quantum platforms continue to scale, this capacity to amortize an expensive training stage into cheap, repeated inference is likely to become increasingly valuable, particularly in feedback-intensive settings such as quantum control~\cite{norambuena2024, Ansel_2024, Dalgaard_2022} and adaptive measurement protocols.

\section{Future work}
\label{sec:conc-future}

The two frameworks presented here open several concrete directions for further research, both as direct extensions and as part of the broader program of applying deep learning to quantum information.

A natural extension of the interpolation framework is to move from reproducing the dynamics to explicitly reconstructing the Hamiltonian. In its current form, the effective Hamiltonian learned for the 7- and 8-qubit systems captures the dynamics without coinciding with the true generator. Incorporating differential constraints derived from the Schrödinger or von Neumann equation into the loss function, in the spirit of the original PINN formulation~\cite{Raissi2019}, could provide a pathway toward genuine Hamiltonian learning while retaining the flexibility of the data-driven approach. This would connect the present work more closely with existing Hamiltonian-reconstruction methods~\cite{che_2021, han_2021, franceschetto2025} and with operator-learning approaches that approximate the propagator directly~\cite{shah2025machine, wang2023machine}. A second promising direction is the extension to open quantum systems~\cite{BreuerPetruccione2006}, where the unitary constraint is replaced by complete positivity and trace preservation, and where the learned dynamical map could be combined with continuous-measurement protocols to further reduce experimental overhead. Finally, architectures capable of generalizing across multiple system sizes or across families of Hamiltonians, rather than requiring a dedicated model per configuration, would substantially broaden the practical applicability of the method, especially for systems with strong time dependence or chaotic behavior where larger architectures are expected to be necessary.

The marginal-reconstruction framework also admits several extensions. The transfer-learning behavior observed across system sizes invites a more systematic study of what the encoder learns, potentially through knowledge-extraction and explainability techniques~\cite{odense2020}, which could yield theoretical insight into the structure of the compatibility conditions themselves. The model's output could also serve as a warm-start initial guess for semidefinite programming solvers~\cite{angell2024, skrzypczyk2023sdp}, combining the speed of learned inference with the exactness of optimization-based methods. Furthermore, the autoencoder architecture is naturally suited to anomaly detection~\cite{islam2020}, suggesting that a related model could learn to distinguish compatible marginals from incompatible ones, addressing the decision version of the Quantum Marginal Problem rather than only its constructive counterpart.

More broadly, the methods developed in this thesis fit within a rapidly expanding effort to apply machine learning across quantum information science. Neural networks have already been used for quantum state reconstruction~\cite{danaci2021, barzili2021}, for solving the quantum many-body problem~\cite{carleo2017, carleo2019, huang2022provably, requena2023}, for learning $N$-representability conditions in quantum chemistry~\cite{sagersmith2022, delgadogranados2025, shao2023}, and for quantum control through physics-informed networks~\cite{norambuena2024}. The two contributions of this thesis, one operating on the level of quantum dynamics and the other on the level of quantum states, complement these efforts and reinforce a consistent conclusion: when the known physics of a problem is built into the learning process, deep learning becomes not a replacement for physical understanding but an extension of it, capable of tackling problems whose exact treatment lies beyond the reach of conventional computation. As both quantum hardware and machine-learning methods continue to mature, the synergy between them explored throughout this thesis is likely to remain a fertile ground for advancing our understanding and control of complex quantum systems.
